\section{Positive kernels and the scalar row estimate}\label{sec:block-recursion}
\hypertarget{source-block-recursion}{}

For a fixed pair of histories the global kernel is a product over places. The Schur product preserves positivity in this product. We can therefore combine the histories in one positive form and apply Cauchy--Schwarz after the local differentiations. Lemma~\ref{lem:residual-single-cofactors} gives the action and relative phase of the two single cofactors in the same compact frame. Poisson summation and a period bound using Fej\'er
kernels turn the resulting Fourier majorants into a scalar
recursion, which is used in the selected-row Gram estimate. Proposition~\ref{prop:positive-block-fourier} specifies the mixed pairing and its two diagonal energies. Their successive suprema are the scalars in \eqref{eq:block-scalar-recurrence}.

\subsection{Fourier bounds from positive blocks}

A derivative can improve the dependence on a single frequency without improving a diagonal energy. Keeping the two diagonal blocks in the positive form lets the mixed estimate use both pieces of information.

\begin{proposition}[Bounds from positive blocks tested on histories]\label{prop:positive-block-fourier}
Fix a positive kernel on the word histories and take the tensor product of the local wave blocks. On a realized marked row, Lemma~\ref{hyp:selected-orbital-coefficient-bounds} identifies the differentiated blocks on those same words, compact types and coefficient frames. The resulting positive form has two diagonal blocks (the two energies) and one off-diagonal block (the mixed pairing).
Let \(l_v\) be the local Fourier frequency and \(X_v>0\) the scale in the fifth-derivative bound at \(v\). Write \(W(x,y)\) for the mixed pairing and \(A'(x,x),B'(y,y)\) for the diagonal forms averaged over the modified words. Then 
\[
 |W(x,y)|\le \frac{A'(x,x)+B'(y,y)}{2}\prod_v\min\{1,(X_v/|l_v|)^5\}.
\]
 A zero frequency contributes a factor of one. The lift comparison has the corresponding majorant in \(t_v^{-1}-t_v'^{-1}\)
with threshold \(2\Lambda'_v\). The weights of the averages
are independent of the vectors testing the forms.
\end{proposition}

\begin{proof}
For a positive block with diagonal forms \(A,B\) and cross form
\(W\), one has
\[
 |x^*Wy|^2\le (x^*Ax)(y^*By).
\]
Testing a positive block on \((x,\lambda y)\) and minimizing
over the phase and modulus of \(\lambda\) proves the displayed
inequality, including when a diagonal value is zero. Conversely,
that inequality bounds the mixed term of the quadratic form by
the sum of the two diagonal terms for every \((x,y)\).

At \(v\in S=\{v:|l_v|>X_v\}\) take the differentiated block. Its entries, with \(|\eta_v|=1\), are
\[
 K^{(v)}_{12}=\eta_v l_v^5 W_v,\qquad
 K^{(v)}_{11}=X_v^5 A'_v,\qquad K^{(v)}_{22}=X_v^5 B'_v ,
\]
where \(A'_v,B'_v\) are positive averages of modified forms. At the other places keep the original block.
Here is why taking the product over places preserves the required inequality for arbitrary histories. On a finite set of indices, including the two block labels, factor the history kernel as \(H_{ij}=\langle h_i,h_j\rangle\). Factor each chosen local block as \(K^{(v)}_{ij}=\langle k_{v,i},k_{v,j}\rangle\), using its positive square root.

The entrywise product is the Gram matrix of explicit vectors:
\[
 H_{ij}\prod_v K^{(v)}_{ij}
   =\left\langle h_i\otimes\bigotimes_v k_{v,i},
                   h_j\otimes\bigotimes_v k_{v,j}\right\rangle .
\]
Thus the mixed entry tested by two arbitrary coefficient families is an inner product of sums of these vectors. Factoring out the powers of \(X_v\) from their squared norms leaves positive averages, with coefficients determined before the vectors are tested. Apply the block inequality to the resulting diagonals \(A',B'\),
with their \(X_v^5\) factors extracted. This yields
\[
 |W(x,y)|\le
  \prod_{v\in S}\left|\frac{X_v}{l_v}\right|^5
                \sqrt{A'(x,x)B'(y,y)}
 \le \prod_v\min\{1,(X_v/|l_v|)^5\}
              \frac{A'(x,x)+B'(y,y)}{2}.
\]
Here each of \(A'\) and \(B'\) is an average in the chosen
modified families. All weights arise from the forms themselves. The zero Fourier coefficient
is the lower-level form by the local integral identity.

For a lifted form the same argument uses the difference of
inverse parameters \(t_v^{-1}-t_v'^{-1}\). At the top of the
chain, if its modulus is greater than \(2\Lambda'_v\), put \(l_v=t_v^{-1}-(t'_v)^{-1}\). The triangle inequality
gives
\[
 \max(|t_v|^{-1},|t'_v|^{-1})\ge |l_v|/2,
 \qquad \min(|t_v|,|t'_v|)\le 2/|l_v|.
\]
At the smaller argument the diagonal has decay exponent ten.
Taking its square root in the block inequality gives exponent
five with threshold \(2\Lambda'_v\). This leaves
three choices per complex place, according to which diagonal
is changed or whether the original block applies. All probability
weights are chosen on the forms before evaluating on the test
vectors. The resulting inequalities can be iterated on the
same arbitrary positive history kernel. \label{cor:block-recursion-1}

\end{proof}

\begin{construction}
\begin{tabular}{@{}lll@{}}\toprule
Frequency at \(v\) & Local block used & Factor in the cross bound\\\midrule
\(|l_v|\le X_v\) & original wave block & \(1\)\\
\(|l_v|>X_v\) & differentiated wave block & \((X_v/|l_v|)^5\)\\\bottomrule
\end{tabular}
\caption{Choice of a local block before taking its Schur product with the history kernel. The local gains multiply, and the averaged diagonal forms can then be tested on the same history.}\label{fig:frequency-block-choice}
\end{construction}

Fourier decay bounds a period by its diagonal energies. Passing to the limiting orbital germ requires an inequality that survives the limit in those energies.

\begin{lemma}[An upper bound for a germ period sum]\label{lem:germ-period}
Use the pairing \(2\RePart\sum_v \xi_v z_v\) on \(\C^f\).
Let \(L_0\) be a full period lattice and \(L_0^*\) its dual for
this pairing. Let \(\phi\) be a nonnegative integrable function,
continuous at the points in a finite set \(S\subset z_0+L_0\).
Assume that a positive block on inverse coordinates has diagonal
corners \(U,Q\) and cross corner \(P\). For any finite
\(T\subset L_0^*\), evaluate it at inverse coordinates
\(x+g\xi\) for \(\xi\in T\). Write \(x=(x_v)_v\in\mathbb C^f\).
Suppose that, with \(T\) fixed and \(\min_v|x_v|\to\infty\), the entries
of \(P\) and \(Q\) converge to
\(\widehat\phi(\xi'-\xi)/\mathcal D\), with \(\mathcal D>0\),
and
\[
 |U(x+g\xi,x+g\xi')|\le M(\xi'-\xi),
 \qquad \sum_{\xi\in L_0^*}M(\xi)<\infty .
\]
Then
\[
 \sum_{z\in S}\phi(z)
 \le \frac{\mathcal D}{\operatorname{covol}(L_0)}
     \sum_{\xi\in L_0^*}M(\xi).
\]
\end{lemma}

\begin{proof}
Test the positive block on \((a,-a)\), with
\(|a_\xi|=|T|^{-1/2}\). This gives
\[
 2\RePart\langle a,Pa\rangle-\langle a,Qa\rangle
 \le \langle a,Ua\rangle
 \le \frac1{|T|}\sum_{\xi,\xi'\in T}M(\xi'-\xi)
 \le \sum_{\xi\in L_0^*}M(\xi).
\]
For the last inequality each difference occurs at most \(|T|\)
times. Pass to the limit
with \(T\) fixed.

Choose a basis \(\mu_1,\ldots,\mu_{2f}\) of \(L_0^*\), and let \(T\) consist of its combinations with coefficients \(0,\ldots,N-1\). Write \([\xi,z]=2\RePart\sum_v\xi_vz_v\). In the Fourier convention \(\widehat\phi(\xi)=\int\phi(z)e^{-2\pi i[\xi,z]}\,dz\), take \(a_\xi=|T|^{-1/2}e^{2\pi i[\xi,z_0]}\). Expanding the square shows that the limiting quadratic form is \(\mathcal D^{-1}\int\phi(z)F_N(z-z_0)\,dz\), where
\[
 F_N(z)=\prod_{j=1}^{2f}\frac1N
       \left|\sum_{h=0}^{N-1}e^{2\pi i h[\mu_j,z]}\right|^2 .
\]
The variables \([\mu_j,z]\bmod 1\) identify normalized Haar measure with product Lebesgue measure on the unit torus.

By character orthogonality each factor has integral one. At distance at least \(\eta\) from an integer, the geometric series bounds that factor by \(1/(N\sin^2(\pi\eta))\), for \(0<\eta<1/2\). Thus the mass outside the coordinate box of radius \(\eta\) tends to zero, while the total unnormalized mass is \(\operatorname{covol}(L_0)\).
At each point of \(S\), integrate over a small neighborhood
where \(\phi\) is bounded below by its value minus \(\delta\).
The neighborhoods can be chosen disjoint, and the Fej\'er
mass in each tends to the full mass of a domain. Summing
these lower bounds and sending \(\delta\) to zero proves
the estimate.
\end{proof}

\subsection{Passing from the germ to a full period}

To pass from one block size to the next, the bound on a germ must control the entire higher period. Write \(U\) for the preceding block and \(\phi\) for the higher diagonal block tested against the positive history kernel. The local Gram blocks give the semidefinite input of
Proposition~\ref{prop:positive-block-fourier}. To use
Lemma~\ref{lem:germ-period}, pair the tensor products at the
same indices with the fixed positive history kernel. The
product of their entries is semidefinite by the Schur product
theorem. This gives the block in that lemma. The Fourier normalizer
\(\mathcal D\) equals the unscaled period covolume
\((d/2)^f g^{3f}d^{2f\nu}\). The lift bound controls \(U\)
at a difference \(g\xi\) by one product majorant times a
positive average of preceding diagonal energies.

\begin{proposition}[Full period through the auxiliary truncation]\label{prop:auxiliary-period-comparison}
Use the history kernel of Proposition~\ref{prop:positive-block-fourier},
with \(R\rho_v/g\ge1\) at the intermediate sizes and
\(\operatorname{vol}(b,s)\le1\) on the actual lattice at the
last size \(s\). In the bound for the full period the truncation is the original one. Its factor in \eqref{eq:block-last-period} is 
\[
 C_4^{fs}\prod_{n=1}^{s-1}\operatorname{vol}(b,n).
\]

\end{proposition}

\begin{proof}
\proofstep{The auxiliary period lattice.}

Suppose these energies are bounded at every point by a
scalar \(P_{n-1}\) times positive averages of the starting
energies. Take first the auxiliary period lattice
\(R\mathbb Z[i]^f\). The sum of the majorant with coordinate
thresholds \(\rho_v>0\) is at most
\[
 C^f\prod_v\bigl(1+(R\rho_v/g)^2\bigr).
\]
For one coordinate, the integer points in a disk of radius
\(T\ge1\) number \(O(T^2)\), and on each successive annulus
of doubled radius the weight falls with exponent five.
This gives the displayed estimate, including thresholds
below one.

Choose \(R\) so that \(R\rho_v/g\ge1\) at every
place. After division by \(R^{2f}\), Lemma~\ref{lem:germ-period}
applied to a singleton gives
\[
 P_n\le C'^f\operatorname{vol}(b,n)P_{n-1},\qquad
 \operatorname{vol}(b,n)=
 \frac{(\prod_v\rho_v)^2}{N(I_D)} .
\]
In this step
\(\mathcal D\prod_v(\rho_v/g)^2
  =(2gd)^{-f}\operatorname{vol}(b,n)\).
The real Euclidean and trace dual conventions change the
absolute constant by a fixed factor per place. Positive
averages are combined by multiplying their coefficients, so
the scalar estimate retains a probability on words. The
initial pointwise estimate follows by the same auxiliary
lattice argument using the wave bound.

\proofstep{The actual period lattice at the last size.}

The auxiliary lattice permits a uniform product estimate during induction. At the last size the normalization belongs to the actual period, whose covolume now cancels the Fourier factor.


Apply the lemma there at size \(s\), with the covolume factor \(\mathcal D\). The dual
frequencies after multiplication by \(g\) lie in \(I_D\).
If \(\operatorname{vol}(b,s)\le1\), Lemma~\ref{lem:lattice-frequency}
with \(H=1\) bounds their total majorant mass by \(C''^f\),
including the zero term. The final coefficient is therefore
\begin{equation}\label{eq:block-last-period}
 C_4^{fs}\prod_{n=1}^{s-1}\operatorname{vol}(b,n).
\end{equation}
Each induction step permits at most five additional letters
at a place. The last volume is used to bound the majorant
mass, and the product of volumes is over the preceding
sizes. The estimate is uniform in the number of places and
their different coordinate thresholds.
\label{cor:block-recursion-2}

\end{proof}

\subsection{The scalar recursion}

The arithmetic application of the recursion uses a different parameter
at each real place. We calculate their product before taking any
supremum over rows. In this calculation \(r=m/2-1\) is the length of
the isotropic flag. Write \(b=e^L\), \(d=|D_E|^{1/(2f)}\),
\(\kappa=2\pi/d\), \(g=g_0g_1\), \(N=10m\), and
\(I=(\Delta-5)/2\). Here \(I\) denotes a real interval length.
For the cut used for two distinct words choose a position \(l\)
at which their bits differ. This is the same position at all real
places by Lemma~\ref{lem:window-parameter-choice}.

\begin{auxlemma}[The frequency scale on two selected words]
\label{lem:selected-frequency-size}
For each block starting at \(i\ge1\) of length \(h\ge1\) with
\(i+h\le r+2\), let \(\rho_{i,h,v}\) be the frequency threshold
at the place \(v\) in the trace coordinate of
Lemma~\ref{lem:period-lattice-normalization}. For the terminal
inverse-coordinate comparison denote the threshold before doubling by
\(\rho'_{i,v}\). Put \(\nu=2r+3-2i-h\). With \(\Theta=k\), for all
sufficiently late selected rows one can take the thresholds to satisfy
\[
 \prod_v(gd^{\nu+1}\rho_{i,h,v})\le\Theta^{-f},\qquad
 \prod_v(gd^{\nu}\rho'_{i,v})\le\Theta^{-f}
       \quad(i+h=r+2).
\]
\end{auxlemma}
\begin{proof}
\proofstep{The cofactor on the interpolation interval.}
Use the row parameters \(\lambda\) and the two column parameters
\(u,u'\) of Definition~\ref{hyp:critical-level-weights} at one place.
Let \(T,T'\) be their assigned primes. The positive size of the
normalized cofactor is given by
\[
 \begin{aligned}
 P(Y)&=\prod_a(Y+\lambda_a),\\
 Q(Y)&=-TT'\prod_{a=1}^k
                   (Y+\tfrac12+u_a)(Y+\tfrac12+u'_a),\\
 B(Y)&=g^m\kappa^{1-m}\sqrt{|P(Y)P(Y+1)/Q(Y)|}.
 \end{aligned}
\]
The numerator is the product of the two shifted ambient central
polynomials. Each compact cofactor contributes one of the two compact
polynomials in the denominator, as in
Lemma~\ref{hyp:selected-orbital-coefficient-bounds}. Write \(z=v_l+Y\).
Use \(5/2\le z\le\Delta/2\) and the shorter interval
\(5/2\le z\le N^3\) for the first singleton. Their suprema of \(B(Y)\)
are denoted by \(B_v,B_v^{(1)}\), respectively.

The negative row coordinates contribute at most
\(\prod_j b_j\exp(4mk\Delta/b_1)\).
For the positive coordinate at position \(a\), put \(w=v_a+Y\).
Except at a replaced endpoint its factor in the square root is
\[
 \sqrt{\frac{|w(w+1)|}
 {|(w-d_a\Delta)(w-d'_a\Delta)|}}.
\]
At the differing bit this factor is
\[
    \sqrt{(z+1)/(\Delta-z)}\le\sqrt{(14/5)z/\Delta}.
\]


For \(a>l\) the factor is at most one. For \(a<l\), write
\(j=l-a\). The factor is at most
\(1+2/(4j-1)\), since \(w\ge4j\Delta\).
For the product use
\[
 \log\prod_{j<k}\left(1+\frac{2}{4j-1}\right)
 \le \sum_{j<k}\left(\frac1{2j}+O(j^{-2})\right)
       \le \tfrac12\log k+O(1).
\]


The first endpoint costs at most \(1/(a_0-v_1)\), where \(a_0\)
is the first definite coordinate. At a top place the replaced
last factor costs at most \(4k\Delta/b_*\) up to an absolute
factor. Indeed its numerator is \(O(k\Delta)\) and its denominator
is at least \(b_*/2\) on the cut. In both cases the endpoint
contribution is \(O(1/a_{\rm eff})\).

Let \(t_v^2=T_{{\rm crit},u,v}T_{{\rm crit},u',v}/(TT')\).
The critical-level calculation in Proposition~\ref{prop:critical-level-lower}
retains the lower bound \(g^m\kappa^{1-m}\prod b_j/(8a_{\rm eff})\)
for either word, with \(4mk\Delta/b_1\le1/10\). It follows that
\[
 B_v^{(1)}\ll t_v\sqrt{k N^3/\Delta},\qquad B_v\ll t_v\sqrt k .
\]
It is the product of the \(t_v\) that is bounded. Let \(Z\) be the
sum of the logarithms of the Euler and finite-index corrections in
Proposition~\ref{prop:prime-target-size}. Its displayed mass formula,
compared with the defining product of the local critical levels, gives
\[
 \sum_v\log(T_{{\rm crit},u,v}/T)
 =f\log(X_u/T)-Z/k
 \le 2\log\Gamma/k+O(fm/k)=O(1).
\]
Here \(T\ge X_u\Gamma^{-2/(fk)}\) and
\(8 C_\Gamma m^5 f L\le k\) by Lemma~\ref{lem:window-parameter-choice}.
Taking the arithmetic mean of the two log estimates gives
\(\prod_v t_v\ll1\).

\proofstep{Frequencies in the trace coordinate.}
For a singleton put \(d_0=2(i-1)\), and for a longer block put
\(d_0=2(i-1)+h-1\). These are the degrees of the leading wave
and lift coefficients in \eqref{eq:word-leading-wave} and
\eqref{eq:word-leading-lift}, respectively. If \(P\) has degree
at most \(d_0>0\), interpolation at the extrema of the rescaled
Chebyshev polynomial on an interval of length \(a\) gives
\[
 |[Y^{d_0}]P|\le \tfrac12(4/a)^{d_0}
                       \max_{Y}|P(Y)|.
\]


Indeed the leading-coefficient interpolation weights alternate in
sign. Their absolute sum is the leading coefficient of the rescaled
Chebyshev polynomial, which takes alternating values \(\pm1\)
at the nodes. For degree zero use evaluation at one point.
In a product of at most \(N\) letters use disjoint translated
subintervals of length at least \(I/(2N)\). The unit shifts in the
consecutive cofactor identities fit in the gaps since
\(I\gg N^3\); the degree-zero first evaluations fit in the short cut.
Thus one degree of interpolation costs at most \(C N/I\), with
\(C\) absolute. The additive character in the unscaled frame
contributes \(2\pi/g\); after reinstating the cofactor, its scale
for an ordinary block is
\[
 \rho_{i,h,v}
 =C(2\pi/g)^{1-m} A_v
              \bigl(C(2\pi/g)N/I\bigr)^{d_0},\qquad
 A_v=\frac{B_v}{g^m\kappa^{1-m}} ,
\]
with \(B_v^{(1)}\) in place of \(B_v\) at \((1,1)\).


On the terminal block \(i+h=r+2\), the one-sided leading term is proportional to \(\tau^{-1}\) on a single compact type. Its interpolation degree is \(i-1\), by the first-component calculation before \eqref{eq:word-leading-lift}. To see the level dependence, put \(x_\rho(Y+1)=\prod_{a<k}(Y+1+\rho_a)\) for the interlacing type in \eqref{eq:word-type-interlacing}. At a node \(Y\) in the middle third of the cut, comparison of the factors on either side of the differing bit gives
\[
  |x_\rho(Y+1)|^2\le \frac{|Q(Y)|}{TT'},\qquad
  |P(Y)|^2\le C A_v^2|Q(Y)|.
\]


For the first inequality each interlacing factor is bounded by the adjoining factor of each of the two words on its side of the cut. The omitted pair at the cut has product of moduli at least one. For the second, every root of \(P\) is outside the cut. For these nodes the definition of \(A_v\) applies since
\[
      |P(Y)/P(Y+1)|\le(1+O(1/I))^{k+m-1}=O(1).
\]

The terminal positive weight for adding a letter is the determinant factor divided by the signed interlacing polynomial at \(Y+1\), as in the conjugated form in Proposition~\ref{prop:real-node-positivity}. For an existing word these factors acquire rational corrections. Each shifted linear factor changes by \(1+O(N^{-2})\) on nodes separated from all earlier nodes and their unit shifts by at least \(N^2\). There are \(O(N)\) factors and \(I\gg N^3\) permits this separation on the same interpolation grids. Hence the ratio of the ambient coefficient to the terminal positive weight is bounded by
\[
        C\frac{|P(Y)|\,|x_\rho(Y+1)|}{|Q(Y)|}
                         \le C A_v/\sqrt{TT'} .
\]

Write \(t\) for the unscaled terminal argument. The frame dilation gives \(|\tau|=(2\pi/g)^{-\nu}|t|\). Taking the squared norm of the one-sided interpolation contributes the square of the sum of the absolute interpolation weights, at most \((CN/I)^{2(i-1)}\) per added letter. Thus the diagonal estimate costs \(|t|^2(\rho'_{i,v})^2\) per letter, where
\[
 (\rho'_{i,v})^2
  =C(2\pi/g)^{-2\nu}(C N/I)^{2(i-1)}A_v/\sqrt{TT'},
 \qquad -2\nu=2i-m .
\]
Five letters give exponent ten on the diagonal. The square root in the positive block inequality gives exponent five for the difference of two inverse coordinates, with threshold \(2\rho'_{i,v}\) by the triangle inequality.

Put \(z_0=C\kappa N/(gI)\). Cancellation of the different and level
powers, using \(d\kappa=2\pi\), gives
\[
 \begin{aligned}
 gd^{\nu+1}\rho_{i,h,v}&=CB_v z_0^{d_0},\\
 (gd^{\nu}\rho'_{i,v})^2&=CB_v\kappa z_0^{2(i-1)}/\sqrt{TT'}
                     \quad(i+h=r+2).
 \end{aligned}
\]
The first line again uses the shorter cut at \((1,1)\).


On the realized family \(m\le f\le k\), \(\kappa\le\pi\),
and \(\sqrt{TT'}\ge b^{m/2}\) by Proposition~\ref{prop:assigned-prime}.
Take the polynomial spacing threshold \(\Delta\ge(kfL)^{12}\)
in Paragraph~\ref{par:window-successive-bounds}; the subsequent
weight width absorbs this polynomial choice as in
Lemma~\ref{hyp:height-section-weights}. Then \(N^3\le\Delta/2\),
\(z_0\ll k/\Delta\), and the four products of normalized radii
are bounded, up to a common constant, by
\[
 \left(C'\sqrt{kN^3/\Delta}\right)^f,\quad
 \left(C'\sqrt k\,z_0^2\right)^f,\quad
 \left(C'\sqrt k\,z_0\right)^f,\quad
 \left(C' k^{1/4}(TT')^{-1/4}\right)^f .
\]


The second bound applies to later singletons, the third to
\(h\ge2\), and the last to the terminal lift. Each base is at
most \(k^{-2}\) eventually, since \(k\) is polynomial in \(x\)
and \(b=\exp(x^\epsilon)\). The factor \(\prod_v t_v\), or its
square root, is bounded independently of \(f\). Increasing the
threshold once proves the assertion with \(\Theta=k\).
\end{proof}


The dependence on a history has now been absorbed into a supremum of positive energies. What remains is to propagate the loss over the whole sequence without multiplying by a large factor at each level.

\begin{proposition}[Accumulated density loss]\label{prop:accumulated-density-loss}
Suppose that \(\sigma_i\) is the square root of the supremum
of the normalized energy at level \(i\), and that the one-level
and long-block comparison satisfies \eqref{eq:block-scalar-recurrence} with
\(0\le a\le q\le 1/2\) and \(e_i\ge0\).
If \((r+1)q\le1\), then at each level \(0\le i\le r+1\) 
\[
 \sigma_i^2\le\bigl(1+C(r+1)q\bigr)\,\sigma_0^2\prod_{j=1}^i(1+e_j),\qquad C=e^6-1.
\]

\end{proposition}

\begin{proof}
In the scalar recursion, write \(\sigma_i\) for the square
root of the supremum of the normalized energy at level \(i\).
Write \(e_i\) for the primitive density errors and take
\(q=(C_0/\Theta)^f\). The terms indexed by \(s\ge2\) use
the bound for the longer blocks.
For the normalized energies the recursion to be estimated is
\begin{equation}\label{eq:block-scalar-recurrence}
 \sigma_i\le (1+a)\sqrt{1+e_i}\,\sigma_{i-1}
       +\sum_{s=2}^i q^{s-1}\sigma_{i-s},
 \qquad 0\le a\le q,\quad e_i\ge0 .
\end{equation}
Put \(D_i=\prod_{j=1}^i\sqrt{1+e_j}\), \(D_0=1\),
and \(a_i=\sigma_i/D_i\). Then \(D_{i-1}\sqrt{1+e_i}=D_i\)
and \(D_{i-s}/D_i\le1\). Hence
\[
 a_i\le(1+a)a_{i-1}+\sum_{s=2}^i q^{s-1}a_{i-s}.
\]

For \(0\le q\le1/2\), the geometric sum is at most \(2q\).
If \(a_j\le \sigma_0(1+3q)^j\) for \(j<i\), every term
\(a_{i-s}\) is bounded by \(\sigma_0(1+3q)^{i-1}\).
Multiplying this bound by \(1+q+2q\) proves the induction.
For \((r+1)q\le1\) the squared factor is bounded by
\[
 (1+3q)^{2(r+1)}\le e^{6(r+1)q}
 \le 1+(e^6-1)(r+1)q .
\]
The last inequality uses the chord bound for the exponential
on \(0\le(r+1)q\le1\). Thus the long-block
errors do not require a constant depending on the number
of recursion levels. The remaining row mass normalization
uses the stabilizer volumes at the selected pair of words.
\end{proof}

\begin{auxlemma}[The accumulated error on the selected row]
\label{lem:selected-density-error}
For two distinct selected words, all starting indices \(i\) and all
permitted block lengths \(h\), in the notation of
Lemma~\ref{lem:selected-frequency-size}, one has
\[
 \operatorname{vol}(i,h)\le k^{-2f},\qquad
 (e^6-1)(r+1)(C_0/k)^f\le b^{-7}.
\]
Here \(C_0\) is an absolute constant that bounds the losses in
\eqref{eq:block-scalar-recurrence} and is fixed before \(x\).
\end{auxlemma}
\begin{proof}
Since \(N(\mathfrak d)=d^{2f}\), the containing dual ideal for this
block has \(N(I_D)=(gd^{\nu+1})^{-2f}\). Thus its volume
is exactly the square of the first product in
Lemma~\ref{lem:selected-frequency-size}. The same calculation for
\(I_L\) uses \(gd^\nu\). The final inverse difference doubles the
radius, so Lemma~\ref{lem:lattice-frequency} applies with \(H=k/2\)
there and with \(H=k\) on \(I_D\). The one-level loss, including
the choices of differentiated blocks, is at most
\[
 a\le (C_1 k^{-5/2})^f\le q,\qquad q=(C_0/k)^f.
\]


For length \(h\ge2\), the primitive density and the one-level
lattice correction each cost at most two on the energy. For the
first of these bounds use the local Euler product
\(\exp(12f2^{-(k-2r-2)})\le2\) on the selected row; the second
follows from the preceding frequency estimate. The norm is the
square root of the energy. The coefficient of the long block is
therefore bounded using \eqref{eq:block-last-period} by
\[
 2 C_4^{fh/2}\prod_{n=1}^{h-1}\operatorname{vol}(i,n)^{1/2}
 \le 2 C_4^{fh/2}k^{-f(h-1)}\le q^{h-1},
\]
once \(C_0\ge2\max(1,C_4,C_1)\). This uses all earlier sizes
at this start. The last volume alone would give only the bound for
the final frequency sum.

Finally \(f\ge 2x^2\), \(r+1\le k\le x^5\) on the late intervals
of Proposition~\ref{prop:actual-rank-budget}. For \(k\ge C_0^2\),
\[
 \log\big((e^6-1)(r+1)q\big)
 \le \log(e^6-1)+\log k-\tfrac12 f\log k\le-7x^\epsilon .
\]
Apply Proposition~\ref{prop:accumulated-density-loss}. The resulting
factor \(1+b^{-7}\) is uniform in the translate, the two words
and the target character with its fixed conductor and weights.
\end{proof}
